\documentclass[10pt,a4paper]{amsart}
\usepackage[margin=19mm]{geometry}
\usepackage{amsmath,amssymb,mathtools}
\usepackage[scr=rsfs,cal=euler]{mathalfa}
\usepackage{xfrac}
\usepackage[unicode,hidelinks,bookmarks=false]{hyperref}
\newcommand{\ie}{i.e.\ }
\newcommand{\cf}{cf.\ }
\newcommand{\resp}{respectively\ }
\newcommand{\Subsection}{\subsection}
\providecommand{\nobreakdash}{\nobreak\hbox{-}}
\setlength{\emergencystretch}{2em}
\multlinegap0pt
\usepackage{amssymb}
\usepackage{xcolor}
\newtheorem{corollary}{Corollary}
\title{Logarithmic Density of Rank $\geq 1$ and Rank $\geq 2$ Genus-2 Jacobians and Applications to Hyperelliptic Curve Cryptography}
\author{Razvan Barbulescu, Mugurel Barcau \and Vicen\c tiu Pa\c sol, George C. \c Turca\c s}
\date{Source: arXiv:2601.17142v2 (CC BY 4.0)}
\begin{document}
\maketitle

\noindent\emph{Source and licence.} Logarithmic Density of Rank $\geq 1$ and Rank $\geq 2$ Genus-2 Jacobians and Applications to Hyperelliptic Curve Cryptography. Version arXiv:2601.17142v2 (CC BY 4.0), distributed by arXiv under the Creative Commons Attribution 4.0 International licence, http://creativecommons.org/licenses/by/4.0/, as recorded in the arXiv licence metadata for this exact version and shown on the abstract page https://arxiv.org/abs/2601.17142v2. Full licence notice: \texttt{licenses/arxiv-2601.17142-CC-BY-4.0.txt}.\par\smallskip

\noindent\textit{Editorial scope.} Counted unit: the article's corollary cor:3 (source0.tex lines 895--913). The article's complete original proof of it is delivered with it (source0.tex lines 914--933, one \texttt{proof} environment of 20 lines). No mathematics is written, repaired or re-derived: the statement and the proof are the article's own lines, in the article's own notation, and no retained line is substituted or re-flowed.  Retained uncounted as the source-local context the proof needs: lines 862--864.  Nothing was omitted from any retained span, so the package carries no omission record.  No retained line contains a citation, so the wrapper carries no bibliography and no bibliography entry.  No retained line contains a figure environment, an image-inclusion command or drawing code, so no image is delivered and the package declares no asset dependency.  Editorial formatting only, in addition: an AMS \texttt{amsart} preamble that restates the article's own declarations and macros used by the retained lines, namely the environments \texttt{corollary} on separate counters, together with the standard packages those lines need.  The retained environments print in this document as Corollary 1 in order of appearance, whereas the article prints the same items with the numbering of its own full document.  The complete native source of the article as distributed in the arXiv e-print for this exact version is bundled unchanged as \texttt{sources/193293/source0.tex}.
\begin{corollary} \label{cor:s1box} The number of pairs $(f,h) \in S_1(X)$ for which $$\left\{ \begin{array}{l} \deg(f) = 6, \\ \Delta(f,h) \neq 0, \\ \mathrm{coeff}_{x^6}(4f(x)+h(x)^2) \text{ is a square, and } \\ \text{the divisor } [\infty_{+} - \infty_{-}] \text{ is torsion on the Jacobian of } y^2+h(x)y=f(x)  \end{array} \right. $$
is of the order $O(X^6(\log X)^6)$, as $X\to \infty$.
\end{corollary}

\begin{corollary} \label{cor:3} For \((f,h)\in S_1^\square(X)\), write \(J_{(f,h)}\)
for the Jacobian of \(y^2+h(x)y=f(x)\). Then, as \(X\to\infty\),
\[
1-O\bigl(X^{-1/2}(\log X)^6\bigr)
\le
\frac{\#\{(f,h)\in S_1^\square(X): \operatorname{rank} J_{(f,h)}(\mathbb Q)\ge 1\}}
{\#S_1^\square(X)}
\le 1.
\]
In particular,
\[
\liminf_{X\to\infty}
\frac{
\log \left|\{(f,h)\in S_1(X): \operatorname{rank} J_{(f,h)}(\mathbb Q)\ge 1\}\right|
}
{\log |S_1(X)|}
\ge \frac{13}{14}.
\]
\end{corollary}

\begin{proof}
Let \(E_1(X)\subset S_1^\square(X)\) be the subset for which
\(\alpha_{(f,h)}=[\infty_+-\infty_-]\in J_{(f,h)}(\mathbb Q)\) is torsion. By
Corollary \ref{cor:s1box} and the count of \(S_1^\square(X)\) above,
\[
\frac{\#E_1(X)}{\#S_1^\square(X)}
=O\bigl(X^{-1/2}(\log X)^6\bigr).
\]
If \(J_{(f,h)}(\mathbb Q)\) has rank \(0\), then every rational point of
\(J_{(f,h)}(\mathbb Q)\) is torsion; in particular \(\alpha_{(f,h)}\) is torsion.
Thus
\[
S_1^\square(X)\setminus E_1(X)
\subseteq
\{(f,h)\in S_1^\square(X): \operatorname{rank}J_{(f,h)}(\mathbb Q)\ge 1\}.
\]
This gives the displayed lower bound, while the upper bound is trivial. Since
\(\#S_1^\square(X)\asymp X^{13/2}\) and \(\#S_1(X)\asymp X^7\), the claimed logarithmic
density lower bound of \(\frac{13}{14}\) also follows.
\end{proof}

\end{document}
